% =============================================================================
%  ANEXO E — ENSAYO CRÍTICO
%  Ciencia, TIC y sociedad contemporánea desde el Nuevo Humanismo
%  -------------------------------------------------------------------------
%  Formato exigido por la consigna:
%    · Letra tamaño 12
%    · Interlineado 1,5
%    · Márgenes estándar (2,54 cm)
%    · Citas y referencias en formato APA
%    · Extensión 3–5 páginas (SIN portada ni referencias)
%  Compilar con:  ./compilar.sh   o   latexmk -pdf main.tex   (usa biber)
% =============================================================================
\documentclass[12pt,letterpaper]{article}

% --- Idioma y codificación ---------------------------------------------------
\usepackage[utf8]{inputenc}
\usepackage[T1]{fontenc}
\usepackage[spanish,es-noshorthands]{babel}
\usepackage{csquotes}

% --- Tipografía (Times 12 pt, estándar APA) ----------------------------------
\usepackage{newtxtext,newtxmath}
\usepackage{microtype}

% --- Márgenes estándar -------------------------------------------------------
\usepackage[letterpaper,margin=2.54cm]{geometry}

% --- Gráficos ----------------------------------------------------------------
\usepackage{graphicx}
\usepackage{xcolor}
\graphicspath{{figs/}}

% --- Encabezado / pie --------------------------------------------------------
\usepackage{fancyhdr}
\pagestyle{fancy}
\fancyhf{}
\fancyhead[R]{\small Ensayo crítico}
\fancyhead[L]{\small\nouppercase{\leftmark}}
\fancyfoot[C]{\thepage}
\renewcommand{\headrulewidth}{0.4pt}

% --- Referencias APA (biblatex + biber) --------------------------------------
\usepackage[backend=biber,style=apa,language=spanish]{biblatex}
\DeclareLanguageMapping{spanish}{spanish-apa}
\addbibresource{referencias/referencias.bib}

% --- Hipervínculos -----------------------------------------------------------
\usepackage[colorlinks=true,linkcolor=black,citecolor=blue,urlcolor=blue]{hyperref}

% --- Interlineado 1,5 --------------------------------------------------------
% Se ajusta con \linespread para evitar el conflicto de setspace con babel
% al parchear \@makecaption. Cambia a 2.0 si requieres interlineado doble.
\linespread{1.5}\selectfont

% =============================================================================
%  DATOS DE LA PORTADA  (editar a conveniencia)
% =============================================================================
\newcommand{\TituloEnsayo}{Título del ensayo crítico}
\newcommand{\SubtituloEnsayo}{Subtítulo aclaratorio (opcional)}
\newcommand{\NombreEstudiante}{Nombre del estudiante}
\newcommand{\CarreraEstudiante}{Ingeniería en Sistemas de Información}
\newcommand{\NombreCurso}{Nombre del curso}
\newcommand{\NombreProfesor}{Nombre del profesor(a)}
\newcommand{\FechaEntrega}{\today}

% =============================================================================
\begin{document}

% --- Portada (NO cuenta para la extensión de 3–5 páginas) --------------------
% =============================================================================
%  Portada — Universidad Nacional
%  No se incluye en la numeración del cuerpo ni en la extensión de 3–5 páginas.
% =============================================================================
\thispagestyle{empty}
\begin{center}
  \vspace*{0.8cm}
  \includegraphics[width=4.4cm]{logo_una.png}
  \vspace{0.9cm}

  {\LARGE\bfseries Universidad Nacional}\\[0.3cm]
  {\large Sede Región Brunca}\\[0.25cm]
  {\large \CarreraEstudiante}

  \vspace{1.6cm}
  {\LARGE\bfseries Los siete pecados capitales en el mundo digital}\\[0.5cm]
  {\Large Cómo la avaricia de las plataformas convierte el vicio en modelo de negocio}
  
  \vspace{1cm}
  \rule{0.65\linewidth}{0.5pt}

  \vspace{0.65cm}
  \renewcommand{\arraystretch}{1.6}
  \setlength{\tabcolsep}{8pt}
  \begin{tabular}{rl}
    \textbf{Curso:}       & CTS en los umbrales del siglo XXI \\
    \textbf{Profesor(a):} & Danny Mata Garbanzo \\
    \textbf{Estudiante:}  & Stward Segura Mendez \\
    \textbf{Carrera:}     & \CarreraEstudiante \\
    \textbf{Fecha:}       & \FechaEntrega \\
  \end{tabular}

  \vspace{0.45cm}
  \rule{0.65\linewidth}{0.5pt}
  \vfill
\end{center}
\newpage


% --- Cuerpo del ensayo -------------------------------------------------------
\setcounter{page}{1}

% =============================================================================
%  INTRODUCCIÓN
% =============================================================================
\section{Introducción}

La Inteligencia Artificial (IA) avanza rápidamente, realizando cada vez más tareas que antes se otorgaban a una persona. Las empresas líderes, movidas por la competitividad y la ambición, buscan maximizar sus ganancias minimizando costos; reemplazar a un ser humano por un sistema que opere sin cansarse y sin parar les es simplemente mejor. Se suele decir que ``la inteligencia artificial es solo una herramienta que aumenta la productividad'', pero esta afirmación, cómoda y tranquilizadora, en un futuro podría convertirse en ``la herramienta que lo hace todo y mejor que todos'': la IA no es un fenómeno reciente, lleva décadas transformando silenciosamente distintos sectores, y su impacto actual no es comparable al de ninguna revolución tecnológica anterior en velocidad ni en alcance.

Como Ingeniero de Sistemas en entrenamiento, esta realidad no me resulta extraña. Los ingenieros somos aquellos encargados de construir estos sistemas, quienes escribimos los algoritmos, quienes diseñamos las arquitecturas que determinarán qué tan rápido y qué tan profundo llegará este cambio. Estamos construyendo algo que tiene el potencial de acabar con el sustento de muchísimas personas ---incluso con el nuestro propio--- y sin embargo, pocas veces en nuestra formación nos detenemos a preguntarnos si deberíamos. El Nuevo Humanismo, entendido en este curso como la búsqueda de un equilibrio entre el progreso tecnológico y el florecimiento genuino del ser humano para ayudar a la sociedad, nos exige exactamente esa pausa. En este ensayo sostengo que el avance de la Inteligencia Artificial, en el contexto de la globalización y el mercado utilitario, profundiza la desigualdad social y amenaza el propósito que muchas personas construyen a través del trabajo; y que la única respuesta coherente desde el Nuevo Humanismo es una ingeniería que se construya desde la consciencia ética, no desde la eficiencia ciega. Quiero recalcar, además, que el propósito y sentido de vida de un ser humano no yace en su utilidad.

% =============================================================================
%  DESARROLLO
% =============================================================================
\section{Desarrollo}

\subsection{La IA no destruye solo empleos: destruye sentido}

Existe una narrativa tranquilizadora que se repite como mantra cada vez que la tecnología desplaza trabajos: ``la revolución industrial también eliminó empleos y al final creó más de los que destruyó''. Esta analogía, sin embargo, ignora una diferencia fundamental. Las revoluciones anteriores mecanizaron la fuerza física; la revolución de la IA mecaniza la cognición. Por primera vez en la historia, la máquina no solo reemplaza lo que nuestros músculos hacen, sino lo que nuestras mentes producen. Y a diferencia del obrero del siglo~XIX que podía reconvertirse en operador de máquina, el trabajador del siglo~XXI desplazado por la IA se enfrenta a la pregunta más desoladora que puede hacerse un ser humano: ¿para qué sirvo?

Los investigadores \textcite{frey2013} estimaron que el 47\,\% de los empleos en Estados Unidos tenían alta probabilidad de automatización en las siguientes dos décadas. Una década después, con la llegada de los modelos de lenguaje de gran escala, esa cifra se antoja conservadora. No se trata solo de empleos rutinarios: se trata de diagnósticos médicos, redacción periodística, análisis legal, generación de código. Se trata, irónicamente, de los trabajos que se consideraban seguros precisamente porque requerían inteligencia.

Lo que más inquieta no es el número de empleos perdidos, sino lo que esos empleos representaban. Para millones de personas, el trabajo no es solo un medio de subsistencia: es identidad, propósito, estructura del tiempo, fuente de pertenencia social. Cuando una empresa automatiza un área, no desaparece únicamente el ingreso; desaparece la razón por la que alguien se levantaba a las cinco de la mañana, el orgullo del programador en el sistema que construyó, la dignidad del contador que resolvía el problema que nadie más podía. La IA no destruye únicamente empleos: destruye, en muchos casos, el sentido que las personas construyeron alrededor de ellos. Y eso, a diferencia de un salario, no se puede reemplazar con un subsidio.

\subsection{La tecnología tiene política: el ingeniero en sistemas como actor social}

Desde la perspectiva de los estudios de Ciencia, Tecnología y Sociedad (CTS), la tecnología nunca es neutral. El politólogo \textcite{winner1980} planteó que los artefactos tecnológicos ``tienen política'': incorporan en su diseño las relaciones de poder, los valores y los intereses de quienes los crean. La Inteligencia Artificial no es la excepción. Cada sistema de IA refleja las prioridades de quienes lo financian, y quienes financian la gran mayoría de estos sistemas son corporaciones cuyo objetivo explícito es la maximización del valor para el accionista. No es una conspiración; es simplemente la lógica del sistema económico en el que operamos.

La globalización intensifica esta dinámica. Los beneficios económicos de la automatización no se distribuyen entre quienes pierden sus empleos, sino que se concentran en las corporaciones transnacionales ---principalmente radicadas en unos pocos polos tecnológicos del mundo--- que desarrollan y comercializan estos sistemas. El \textcite{wef2020} estimó que para 2025 la automatización destruiría 85 millones de puestos de trabajo a nivel global mientras creaba 97 millones nuevos. La diferencia parecería positiva sobre el papel, pero esos 97 millones de empleos nuevos requieren competencias altamente especializadas, acceso a educación de calidad y conectividad: recursos distribuidos de manera profundamente desigual. El resultado no es un mundo con más empleo: es un mundo donde la brecha entre quienes pueden acceder a la economía digital y quienes no, se vuelve abismal.

Como ingeniero en sistemas, esta realidad me coloca en una posición incómoda que no puedo ignorar. Yo soy parte de la maquinaria que produce esta transformación. Cada línea de código que escribo, cada sistema que diseño, cada solución que optimizo forma parte de un proceso que tiene consecuencias sobre personas reales. Esta incomodidad es precisamente la que el Nuevo Humanismo quiere que sintamos, porque es el punto de partida de una reflexión ética genuina. \textcite{han2012} describe nuestra época como una \textit{sociedad del rendimiento}, donde el imperativo de la eficiencia ha reemplazado al imperativo del sentido. Un ingeniero formado únicamente en rendimiento ---en cuánto procesa, cuánto optimiza, cuánto escala--- es exactamente el tipo de profesional que esa sociedad produce y reproduce. Un ingeniero formado también en humanismo es, quizás, el que puede interrumpir el ciclo.

El problema es que nuestra formación académica raramente nos prepara para esa interrupción. Aprendemos estructuras de datos, algoritmos, paradigmas de programación, arquitecturas de software. Aprendemos a resolver problemas técnicos con elegancia y eficiencia. Lo que rara vez aprendemos es a preguntarnos si el problema que estamos resolviendo vale la pena resolverlo, o a quién beneficia realmente la solución. Como señala \textcite{borgmann1984}, la tecnología moderna actúa como un ``dispositivo'' que oculta su mecanismo y entrega únicamente su producto final, separando al usuario ---y al mismo creador--- de la complejidad y del costo real de lo que produce. La IA es, en muchos sentidos, el dispositivo definitivo: genera valor sin mostrar el costo humano, produce eficiencia sin preguntar a quién le cuesta esa eficiencia.

\subsection{El Nuevo Humanismo como respuesta y como exigencia}

El Nuevo Humanismo, tal como lo concibe la \textcite{unesco2015}, no es una nostalgia por un mundo sin tecnología: es una convicción de que el desarrollo humano debe ser el centro y el fin de toda transformación social, incluida la tecnológica. En ese marco, la pregunta no es si debemos desarrollar IA ---esa decisión ya fue tomada por fuerzas que ningún individuo controla--- sino cómo debemos desarrollarla, bajo qué principios, con qué límites y en función de qué valores.

\textcite{harari2016} plantea un escenario tan probable como perturbador: una humanidad dividida entre una pequeña élite de personas relevantes para la economía algorítmica y una vasta mayoría que simplemente ha quedado obsoleta. No se trata de ciencia ficción distante; se trata de una extrapolación razonable de las tendencias actuales. Y lo que hace que este escenario sea tan profundamente triste no es la pobreza material que implicaría, sino la pérdida de propósito. Una persona sin un rol reconocido en la sociedad, sin la experiencia de ser necesaria, no es solo una persona pobre: es una persona sin sentido.

El Nuevo Humanismo responde a esto con una afirmación radical: el valor del ser humano no puede depender de su productividad económica. Esta afirmación, que parece obvia cuando se enuncia, es subversiva dentro del sistema de valores del mercado global. Nos dice que hay que diseñar tecnología no para maximizar la rentabilidad de las corporaciones, sino para ampliar las posibilidades de la vida humana; para distribuir el ocio liberado por la automatización en lugar de concentrar la riqueza que genera; para emplear la eficiencia de las máquinas en resolver los problemas que el mercado nunca resolverá porque no son rentables: el cambio climático en comunidades pobres, el acceso a la salud en regiones marginadas, la educación de calidad para quienes no pueden pagarla.

Este es el reto que tenemos como ingenieros en sistemas: no somos solo técnicos que ejecutan especificaciones. Somos actores sociales con una capacidad real de influir en cómo se despliega esta tecnología. Y esa capacidad viene acompañada de una responsabilidad que no podemos tercerizar a los gerentes, a los inversores o al mercado.

% =============================================================================
%  CONCLUSIONES
% =============================================================================
\section{Conclusiones}

La solución apunta al Nuevo Humanismo: un enfoque que busca el bien común mediante un uso responsable y ético de la tecnología. Muchos ingenieros en sistemas, programadores y diseñadores fueron los responsables de crear las plataformas que hoy causan daño a las nuevas generaciones; precisamente por eso, el cambio debe empezar desde nuestra propia formación. Como futuro ingeniero en sistemas, comprendo que la pregunta no es solo si una tecnología funciona, sino a quién beneficia y a quién perjudica su diseño. Si los pecados capitales fueron incrustados en el código por avaricia, también pueden ser desincrustados por una ética del bien común. Cambiar a los demás es difícil, pero cambiarse a uno mismo lo es aún más; por eso el primer paso es adoptar y practicar el Nuevo Humanismo de manera individual, para luego inculcarlo y generar un cambio positivo. Quizás ese sea el verdadero reto de mi generación de ingenieros: construir un mundo digital que se mueva por el bien común y no por la avaricia.



% --- Referencias (NO cuentan para la extensión de 3–5 páginas) ---------------
\newpage
\printbibliography[title={Referencias}]

\end{document}
